\documentclass[svenska]{style/LTHthesis}

% Preamble
% from template
\usepackage[T1]{fontenc}%
\usepackage{mathptmx, helvet}%
\usepackage{graphicx}%

\usepackage{wrapfig}
\usepackage{url}
\raggedbottom

\usepackage{xcolor}
\usepackage{hyperref}
\hypersetup{
  colorlinks,
  allcolors={blue!40!black},
}

\pagenumbering{gobble}
%
%%%%%%%%%%%%%%%%%%%%%%%%%
%%% Unified notation
%%%%%%%%%%%%%%%%%%%%%%%%%

%% Mathematics
% Shorthand
\newcommand{\R}{\mathbb{R}}
\newcommand{\T}{\mathrm{\small{T}}}

% Functions
\newcommand{\LL}{\mathcal{L}_2}
\newcommand{\funof}[1]{\left( #1 \right)}
\newcommand{\abs}[1]{\left|#1\right|}
\newcommand{\norm}[1]{\lVert#1\rVert}
\newcommand{\E}[1]{\mathbb{E}\left[ #1 \right]}
\newcommand{\trace}[1]{\mathrm{tr}\left(#1\right)}
\newcommand{\diag}[1]{\mathrm{diag}\left\{#1\right\}}

%% Control
\newcommand{\plant}{\mathcal{P}}
\newcommand{\ctrler}{\mathcal{C}}

%% Code

%% Real-Time
%

% Document
\begin{document}

\noindent\begin{minipage}[]{0.25\textwidth}
    \noindent\includegraphics[width=20mm]{style/LUCsv}
    \par\vspace*{3mm} 
\end{minipage}\hfill
\begin{minipage}[]{0.74\textwidth}
    \vspace*{-3mm}
    \sffamily
    \begin{flushleft}
        \huge Ett viktigt bidrag till teorin\\[2mm]
        \Large Nils Vreman\\[1mm]
        \normalsize Institutionen för Reglerteknik
    \end{flushleft}
\end{minipage}
\par\vspace{1mm}  
\begin{sffamily}\noindent Populärvetenskaplig sammanfattning av doktorsavhandlingen \emph{A Very Important Contribution to the Theory}, juni 2023. Avhandlingen kan laddas ner från: \url{http://www.control.lth.se/publications}

\end{sffamily}
\vspace{2mm}\hrule\vspace{3mm}
\noindent
Denna avhandling undersöker hur tillämpningar  kan förbättras på olika vis. I dagens industri används reglerteknik till många olika uppgifter, men till största delen rör det sig om enkla uppdrag som inte är beroende av fysisk kontakt med omgivningen.

\begin{wrapfigure}{r}{0.45\textwidth}
    \vspace*{-2mm}
    \begin{centering}
        \fbox{\rule{0pt}{30mm}\rule{0.95\linewidth}{0pt}}
    \end{centering}\small\par\medskip
    \noindent\emph{En industrirobot utrustad med fräsverktyg, som använts för experiment i avhandlingen.}
    \vspace*{-2mm}
\end{wrapfigure}

För bearbetningsuppgifter såsom slipning och fräsning används traditionellt sett specialmaskiner. Intresset för att istället använda industrirobotar har ökat, på grund av robotars flexibilitet i att kunna utföra en mängd olika uppgifter, samt den relativt låga investeringskostnaden. Industrirobotar är mycket bra på att positionera sig noggrant i fria luften, men under bearbetning när stora krafter utövas på roboten blir resultatet inte lika bra. Detta beror på att robotens leder är förhållandevis veka och böjs när externa krafter påverkar roboten, vilket i sin tur leder till att noggrannheten blir lidande. Om man hade kunnat mäta ledernas utböjning direkt hade det varit enkelt att kompensera för dem, men i dagsläget saknas sensorer som kan mäta utböjningen.

Här skall texten fortsätta, tills den blir antingen två eller (gärna) fyra 
sidor lång.
\end{document}
